% These are SEPARATE replacements/insertions, not one continuous section.
% INSERTION 1: add beside the embedding model in the Stage 3 setup.
% These two distinct citations bring the total to 18.
Semantic similarity is computed using the
\texttt{sentence-transformers/all-mpnet-base-v2} checkpoint, an MPNet-based
sentence embedding model in the Sentence-BERT framework~\cite{reimers2019sentence,song2020mpnet}.
The exact checkpoint is documented in its model card.\footnote{\url{https://huggingface.co/sentence-transformers/all-mpnet-base-v2}}

% REPLACEMENT 2: Stage 2 metric. Preserves the frozen evaluation procedure.
Stage~2 uses the frozen coarse sentence-level overlap view. Within each sample
and semantic field, an extracted span is counted as matched if it overlaps at
least one reference region; a reference region is counted as matched if it
overlaps at least one extracted span. These are separate counts, rather than
a one-to-one matching. Pooling over actor, action, condition, constraint, and
exception, precision and recall are
\begin{equation}
P=\frac{N_{\mathrm{matched\ predictions}}}{N_{\mathrm{predictions}}},\qquad
R=\frac{N_{\mathrm{matched\ references}}}{N_{\mathrm{references}}},
\end{equation}
and pooled F1 is \(2PR/(P+R)\). Modality is evaluated separately using
four-class Macro-F1 and is excluded from this pooled span score.
Non-empty overlap measures partial evidence recovery; it does not establish
exact boundaries or complete semantic correctness.

% REPLACEMENT 3: main-table versus prompt-ablation identity.
% Verify per-run settings separately; do not infer an unrecorded release.
Table~I uses the frozen v6/R3 extraction snapshot registered as
\texttt{deepseek-v4-pro}, whereas Table~II evaluates the separate modular
E/S/J prompt family. The full modular configuration~111 is therefore not the
same prompt snapshot as the Table~I extractor. The eight configurations were
each evaluated once across two execution batches; their differences are
descriptive and do not establish stability across repeated runs. The basic
task and schema requirements remain present in configuration~000.

% REPLACEMENT 4: non-LLM baseline.
Rules-Only is our independent reconstruction of the approach of
Sun et al.~\cite{sun2024design}. It combines a BERT-TextCNN modality classifier
with syntactic analysis and manually specified extraction rules. Thus,
Rules-Only denotes the non-generative-LLM baseline and does not imply that
every component is an untrained rule system. Its scores are results of this
local reconstruction, not scores reported in Sun et al.'s original paper.

% REPLACEMENT 5: compact limitations paragraph after the results.
The findings have several limits. The extraction benchmark contains 150
requirements with reconstructed reference annotations, and the overlap metric
does not require exact span boundaries. The prompt configurations have one
observed run each. The downstream GDPR benchmark was also used during
development and threshold calibration, so its scores do not estimate
performance on an independently held-out test set. Finally, higher extraction
scores do not remove errors in aligning rule semantics with BPMN labels, and
the three evaluated violation types do not demonstrate complete checking of
every semantic element in the Rule Record.
